\section{Analytical Capabilities}


\subsection{Live Dashboards}

Dashboards query materialized views for sub-second response:

\begin{lstlisting}[language=SQL,caption=Dashboard query]
-- Orders and revenue, last hour, by minute
SELECT
    window_start,
    SUM(order_count) AS orders,
    SUM(total_revenue) AS revenue
FROM mv_revenue_by_minute
WHERE store_id = 'store_123'
  AND window_start >= NOW() - INTERVAL '1' HOUR
GROUP BY window_start
ORDER BY window_start;
\end{lstlisting}

Query latency: 15--50ms (served entirely from Redis).

\subsection{Cohort Analysis}

Cohort analysis groups users by acquisition characteristics and tracks behavior over time:

\begin{lstlisting}[language=SQL,caption=Cohort retention query]
WITH cohorts AS (
    SELECT
        user_id,
        DATE_TRUNC('week', MIN(event_time)) AS cohort_week
    FROM events
    WHERE event_type = 'purchase'
      AND store_id = 'store_123'
    GROUP BY user_id
),
activity AS (
    SELECT
        e.user_id,
        c.cohort_week,
        DATE_TRUNC('week', e.event_time) AS activity_week
    FROM events e
    JOIN cohorts c ON e.user_id = c.user_id
    WHERE e.event_type = 'purchase'
      AND e.store_id = 'store_123'
)
SELECT
    cohort_week,
    DATE_DIFF('week', cohort_week, activity_week) AS weeks_since_cohort,
    COUNT(DISTINCT user_id) AS active_users
FROM activity
GROUP BY cohort_week, weeks_since_cohort
ORDER BY cohort_week, weeks_since_cohort;
\end{lstlisting}

We provide cohort analysis as a built-in function:

\begin{lstlisting}[language=SQL,caption=Cohort analysis function]
SELECT * FROM COHORT_ANALYSIS(
    events,
    cohort_date := DATE_TRUNC('week', first_purchase_date),
    activity_date := DATE_TRUNC('week', event_time),
    user_id := user_id,
    metric := COUNT(DISTINCT order_id),
    filters := [store_id = 'store_123', event_type = 'purchase']
);
\end{lstlisting}

\subsection{Funnel Analysis}

Funnel analysis tracks conversion through defined steps:

\begin{lstlisting}[language=SQL,caption=Conversion funnel]
SELECT * FROM FUNNEL(
    events,
    steps := [
        event_type = 'page_view' AND page = '/products',
        event_type = 'product_view',
        event_type = 'add_to_cart',
        event_type = 'checkout_start',
        event_type = 'purchase'
    ],
    window := INTERVAL '30' MINUTE,
    group_by := [device_type, utm_source]
);
\end{lstlisting}

The funnel operator uses session-based windowing to track users through steps:

\begin{lstlisting}[language=Python,caption=Funnel computation]
def compute_funnel(events, steps, window, group_by):
    results = defaultdict(lambda: [0] * len(steps))

    for session in group_by_session(events):
        # Track furthest step reached per session
        step_times = [None] * len(steps)

        for event in sorted(session.events, key=lambda e: e.timestamp):
            for i, step_predicate in enumerate(steps):
                if step_predicate(event):
                    if i == 0 or (step_times[i-1] and
                        event.timestamp - step_times[i-1] <= window):
                        step_times[i] = event.timestamp

        # Increment counts for reached steps
        group_key = tuple(getattr(session, g) for g in group_by)
        for i, t in enumerate(step_times):
            if t is not None:
                results[group_key][i] += 1

    return results
\end{lstlisting}

\subsection{Anomaly Detection}

Verus integrates anomaly detection into the analytics pipeline:

\subsubsection{Statistical Anomalies}

We detect anomalies using seasonal decomposition:

\begin{equation}
y_t = T_t + S_t + R_t
\end{equation}

where $T_t$ is trend, $S_t$ is seasonal component, and $R_t$ is residual. Anomalies are residuals exceeding threshold:

\begin{equation}
|R_t| > k \cdot \sigma_R
\end{equation}

\begin{lstlisting}[language=Python,caption=Anomaly detection]
class AnomalyDetector:
    def __init__(self, sensitivity=3.0):
        self.k = sensitivity

    def detect(self, series: pd.Series) -> List[Anomaly]:
        # Seasonal decomposition
        decomposition = seasonal_decompose(series, period=24*7)  # Weekly
        residual = decomposition.resid.dropna()

        # Detect anomalies
        threshold = self.k * residual.std()
        anomalies = []

        for timestamp, value in residual.items():
            if abs(value) > threshold:
                anomalies.append(Anomaly(
                    timestamp=timestamp,
                    expected=series[timestamp] - value,
                    actual=series[timestamp],
                    severity=abs(value) / threshold
                ))

        return anomalies
\end{lstlisting}

\subsubsection{Alert Integration}

Anomalies trigger alerts through configurable channels:

\begin{lstlisting}[language=SQL,caption=Alert configuration]
CREATE ALERT order_drop AS
SELECT
    store_id,
    COUNT(*) AS order_count,
    AVG(COUNT(*)) OVER (
        PARTITION BY store_id
        ORDER BY window_start
        ROWS BETWEEN 168 PRECEDING AND 1 PRECEDING  -- Last week
    ) AS expected
FROM mv_orders_by_hour
GROUP BY store_id, window_start
HAVING order_count < expected * 0.5  -- 50% drop
WITH NOTIFICATION slack_webhook, email;
\end{lstlisting}

